% =====================================================================
\chapter{Optimización continua y convexa}
\label{tema:4}
% =====================================================================

Cuando las variables son continuas y el problema es convexo, todo mínimo local es global, existen certificados de optimalidad y los multiplicadores tienen interpretación económica. Este tema sigue los capítulos~4 («Convex optimization problems») y~5 («Duality»), en especial las secciones~5.5 («Optimality conditions») y~5.6 («Perturbation and sensitivity analysis»), de Boyd y Vandenberghe~\cite{bv}, texto del curso EE364a de Stanford; el capítulo~12 de Nocedal y Wright~\cite{nw} para las condiciones de Karush--Kuhn--Tucker en problemas no convexos; y el capítulo~3 («Economic Dispatch of Thermal Units and Methods of Solution») de Wood, Wollenberg y Sheblé~\cite{wood} para el caso de estudio. Las herramientas son CVXPY~\cite{cvxpy,cvxpydoc} y SciPy. El módulo~7 (tema~1) introdujo la convexidad en el contexto del entrenamiento; aquí se añaden las restricciones.

\section{Problemas convexos}

\begin{definicion}[Problema convexo]\label{def:convexo}
El problema de la \cref{def:prob-opt} es \emph{convexo} si $X=\R^{n}$, $f$ y las $g_i$ son funciones convexas y las $h_j$ son afines, $h_j(x)=a_j\T x-b_j$. Su conjunto factible es entonces convexo, como intersección de subniveles de funciones convexas e hiperplanos.
\end{definicion}

\begin{teorema}[Local implica global]\label{teo:local-global}
En un problema convexo, todo óptimo local es óptimo global.
\end{teorema}

\begin{proof}
Sea $x^{*}$ óptimo local en $\mathcal{X}\cap B(x^{*},r)$ y supóngase que existe $\tilde x\in\mathcal{X}$ con $f(\tilde x)<f(x^{*})$. Para $\theta\in(0,1]$ el punto $x_\theta=(1-\theta)x^{*}+\theta\tilde x$ es factible, por la convexidad de $\mathcal{X}$, y $f(x_\theta)\le(1-\theta)f(x^{*})+\theta f(\tilde x)<f(x^{*})$, por la de $f$. Tomando $\theta$ pequeño, $x_\theta\in B(x^{*},r)$, lo que contradice la optimalidad local.
\end{proof}

Las clases convexas más usadas, de menor a mayor generalidad, son la programación lineal, la cuadrática (objetivo $\tfrac12x\T Px+q\T x$ con $P\succeq 0$ y restricciones lineales), la cónica de segundo orden y la semidefinida. CVXPY sigue las reglas de la \emph{programación convexa disciplinada} (DCP): construye las expresiones con funciones atómicas de curvatura conocida y reglas de composición que garantizan la convexidad, de modo que rechaza un modelo cuya convexidad no puede certificar, y lo traduce a la forma cónica que aceptan resolvedores como Clarabel, SCS u OSQP~\cite{cvxpy,cvxpydoc}.

\section{Lagrangiana, dualidad y condiciones KKT}

\begin{definicion}[Lagrangiana y función dual]\label{def:lagrangiana}
Para el problema de la \cref{def:prob-opt} con $X=\R^{n}$, la \emph{lagrangiana} es
\[
\Lag(x,\lambda,\nu)=f(x)+\sum_{i=1}^{m}\lambda_ig_i(x)+\sum_{j=1}^{p}\nu_jh_j(x),\qquad\lambda\ge 0,
\]
y la \emph{función dual} es $q(\lambda,\nu)=\inf_x\Lag(x,\lambda,\nu)$.
\end{definicion}

\begin{proposicion}[Dualidad débil]\label{prop:dual-debil-conv}
Si $\lambda\ge 0$, entonces $q(\lambda,\nu)\le f(x)$ para todo $x$ factible; por tanto $q(\lambda,\nu)\le p^{*}$, el valor óptimo.
\end{proposicion}

\begin{proof}
Si $x$ es factible, $g_i(x)\le 0$ y $h_j(x)=0$, luego $\sum_i\lambda_ig_i(x)+\sum_j\nu_jh_j(x)\le 0$ y $\Lag(x,\lambda,\nu)\le f(x)$. El ínfimo sobre todos los $x$ es menor o igual que el valor en ese $x$.
\end{proof}

En un problema convexo que cumple la \emph{condición de Slater} (existe un punto estrictamente factible, $g_i(x)<0$ para las restricciones no afines), la dualidad es fuerte: $\max q=p^{*}$ (Boyd y Vandenberghe~\cite{bv}, cap.~5).

\begin{teorema}[Condiciones de Karush--Kuhn--Tucker]\label{teo:kkt}
Sean $f$, $g_i$ y $h_j$ diferenciables. Se dice que $(x^{*},\lambda^{*},\nu^{*})$ cumple las condiciones KKT si
\begin{enumerate}
  \item \emph{factibilidad primal}: $g_i(x^{*})\le 0$, $h_j(x^{*})=0$;
  \item \emph{factibilidad dual}: $\lambda^{*}\ge 0$;
  \item \emph{holgura complementaria}: $\lambda_i^{*}g_i(x^{*})=0$ para todo $i$;
  \item \emph{estacionariedad}: $\nabla f(x^{*})+\sum_i\lambda_i^{*}\nabla g_i(x^{*})+\sum_j\nu_j^{*}\nabla h_j(x^{*})=0$.
\end{enumerate}
Si el problema es convexo y se cumplen las condiciones KKT, $x^{*}$ es óptimo global. Si el problema es convexo y cumple la condición de Slater, todo óptimo satisface las KKT para algún $(\lambda^{*},\nu^{*})$~\cite{bv}. En problemas no convexos, las KKT son necesarias en un mínimo local bajo una cualificación de restricciones, pero no suficientes~\cite{nw}.
\end{teorema}

\begin{proof}[Demostración de la suficiencia]
Con $\lambda^{*}\ge 0$ y $h_j$ afines, $x\mapsto\Lag(x,\lambda^{*},\nu^{*})$ es convexa y diferenciable, y por (iv) su gradiente se anula en $x^{*}$; una función convexa alcanza su mínimo donde se anula el gradiente, así que $q(\lambda^{*},\nu^{*})=\Lag(x^{*},\lambda^{*},\nu^{*})$. Por (i) y (iii), $\Lag(x^{*},\lambda^{*},\nu^{*})=f(x^{*})$. Para cualquier $x$ factible, por la \cref{prop:dual-debil-conv}, $f(x)\ge q(\lambda^{*},\nu^{*})=f(x^{*})$.
\end{proof}

\begin{proposicion}[Multiplicadores como sensibilidades]\label{prop:sensibilidad}
Sea el problema convexo perturbado $\min f(x)$ s.a. $g_i(x)\le u_i$, $h_j(x)=w_j$, con valor óptimo $p^{*}(u,w)$, y sea $(\lambda^{*},\nu^{*})$ dual óptimo del problema sin perturbar con dualidad fuerte. Entonces
\[
p^{*}(u,w)\ \ge\ p^{*}(0,0)-\lambda^{*\mathsf{T}}u-\nu^{*\mathsf{T}}w ,
\]
y si $p^{*}$ es diferenciable en $(0,0)$, $\partial p^{*}/\partial u_i=-\lambda_i^{*}$ y $\partial p^{*}/\partial w_j=-\nu_j^{*}$.
\end{proposicion}

\begin{proof}
Para $x$ factible en el problema perturbado, $f(x)\ge f(x)+\sum_i\lambda_i^{*}(g_i(x)-u_i)+\sum_j\nu_j^{*}(h_j(x)-w_j)\ge q(\lambda^{*},\nu^{*})-\lambda^{*\mathsf{T}}u-\nu^{*\mathsf{T}}w$, y $q(\lambda^{*},\nu^{*})=p^{*}(0,0)$. Tomando el ínfimo en $x$ se obtiene la desigualdad. Si $p^{*}$ es diferenciable, la desigualdad dice que el plano $p^{*}(0,0)-\lambda^{*\mathsf{T}}u-\nu^{*\mathsf{T}}w$ queda por debajo de la función y la toca en el origen, de modo que coincide con su plano tangente (Boyd y Vandenberghe~\cite{bv}, §5.6).
\end{proof}

Es la generalización de los precios sombra de la \cref{prop:sombra}: relajar una restricción activa abarata el óptimo, en primera aproximación, en su multiplicador por unidad.

\section{Caso de estudio: despacho económico}

\begin{ejemplo}[Despacho económico de tres grupos térmicos]\label{ej:despacho}
Un operador debe cubrir una demanda de $D=850$~MW durante una hora con tres grupos térmicos. El coste horario de cada grupo, ajustado a sus curvas de consumo, es cuadrático, $C_i(P_i)=a_i+b_iP_i+c_iP_i^{2}$~(€/h):
\begin{center}\small
\begin{tabular}{@{}lccccc@{}}
\toprule
Grupo & $a_i$ (€/h) & $b_i$ (€/MWh) & $c_i$ (€/MW$^2$h) & $P_i^{\min}$ (MW) & $P_i^{\max}$ (MW) \\
\midrule
1 & 500 & 18 & 0,020 & 100 & 400 \\
2 & 400 & 20 & 0,010 & 50 & 350 \\
3 & 300 & 26 & 0,005 & 50 & 300 \\
\bottomrule
\end{tabular}
\end{center}
\emph{Modelo.} $\min\sum_iC_i(P_i)$ s.a. $\sum_iP_i=D$ y $P_i^{\min}\le P_i\le P_i^{\max}$. Como $c_i>0$, el objetivo es convexo y las restricciones son lineales: es un programa cuadrático convexo y se cumple la condición de Slater (por ejemplo, $P=(300,300,250)$ es estrictamente factible en los límites). \emph{Elección del método:} KKT a mano para entender la solución, CVXPY para resolverla y verificarla.

\emph{Condiciones KKT.} Con multiplicador $\lambda$ para el balance y $\mu_i^{\pm}\ge 0$ para los límites, la estacionariedad es $b_i+2c_iP_i-\lambda-\mu_i^{-}+\mu_i^{+}=0$. Si ningún grupo está en un límite, todos operan al mismo \emph{coste incremental}:
\[
b_i+2c_iP_i=\lambda\quad\Longrightarrow\quad P_i=\frac{\lambda-b_i}{2c_i}=25(\lambda-18),\ 50(\lambda-20),\ 100(\lambda-26).
\]
\emph{Paso 1.} Sumando e igualando a $D$: $175\lambda-4050=850$, de donde $\lambda=28$ y $P=(250,400,200)$; pero $P_2=400>350$ viola el límite.
\emph{Paso 2.} Se fija $P_2=350$ (límite activo) y se reparte el resto, $500$~MW, entre los grupos 1 y 3: $25(\lambda-18)+100(\lambda-26)=500$, es decir, $125\lambda=3550$ y $\lambda=28{,}4$~€/MWh. Entonces $P_1=25\cdot 10{,}4=260$ y $P_3=100\cdot 2{,}4=240$~MW, ambos dentro de sus límites.
\emph{Paso 3.} El multiplicador del límite del grupo~2 es $\mu_2^{+}=\lambda-(b_2+2c_2P_2)=28{,}4-27=1{,}4\ge 0$: las KKT se cumplen y, por el \cref{teo:kkt}, la solución es óptima.

\emph{Resultado.} $\bm{P^{*}=(260,\ 350,\ 240)}$~MW, con coste \textbf{21\,985~€/h} ($6532+8625+6828$). El multiplicador $\lambda=28{,}4$~€/MWh es el precio marginal del sistema: por la \cref{prop:sensibilidad}, 10~MW más de demanda costarían unos $284$~€/h; el valor exacto, resolviendo de nuevo, es $284{,}4$~€/h, porque el coste es cuadrático. El multiplicador $\mu_2^{+}=1{,}4$~€/MWh indica cuánto se ahorraría por MW si el grupo~2 pudiera producir más. CVXPY reproduce la solución y devuelve los mismos multiplicadores (cuaderno, parte~4):
\begin{verbatim}
P = cp.Variable(3); bal = cp.sum(P) == D
coste = cp.sum(a + cp.multiply(b, P) + cp.multiply(c, cp.square(P)))
cp.Problem(cp.Minimize(coste), [bal, P >= Pmin, P <= Pmax]).solve()
bal.dual_value   # 28,4 en valor absoluto (signo según convenio)
\end{verbatim}
\end{ejemplo}

\begin{observacion}[Cuando se pierde la convexidad]\label{obs:noconvexo}
El caso anterior deja de ser convexo con detalles realistas: costes con efecto de apertura de válvulas (términos sinusoidales), mínimos técnicos con arranque y parada (variables binarias, el problema de asignación de unidades del cap.~4 de~\cite{wood}) o mezclas con calidades que se multiplican por caudales (términos bilineales). Entonces las KKT son solo necesarias, los métodos locales (SQP, punto interior) pueden quedarse en óptimos locales y conviene arrancar desde varios puntos, usar un resolvedor global o reformular. En el ajuste de redes neuronales del módulo~7 el problema tampoco es convexo, pero allí basta con un buen mínimo local.
\end{observacion}

\begin{table}[h]
\centering\small
\caption{Tecnologías del tema~4: optimización continua.}\label{tab:tec4}
\begin{tabularx}{\textwidth}{@{}l>{\raggedright\arraybackslash}X>{\raggedright\arraybackslash}p{2.2cm}>{\raggedright\arraybackslash}X@{}}
\toprule
Herramienta & Para qué sirve & Licencia o tipo & Cuándo elegirla \\
\midrule
CVXPY~\cite{cvxpy} & modelización convexa con verificación DCP y duales & abierta (Apache~2.0) & cualquier problema convexo de tamaño medio \\
Clarabel, SCS, OSQP & resolvedores cónicos (Clarabel, SCS) y cuadráticos (OSQP) & abiertos (Apache~2.0; SCS, MIT) & los usa CVXPY; OSQP, para control predictivo rápido \\
SciPy \texttt{minimize} & SLSQP, \texttt{trust-constr}, L-BFGS-B & abierta (BSD-3) & no lineal local con pocas variables \\
Ipopt & punto interior no lineal a gran escala & abierta (EPL-2.0) & no lineal grande con derivadas (vía Pyomo o CasADi) \\
Gurobi & cuadrático y no convexo bilineal global & comercial & mezclas y redes con términos bilineales \\
\bottomrule
\end{tabularx}
\end{table}

\begin{observacion}[Errores frecuentes en optimización continua]\label{obs:errores4}
\begin{enumerate}[label=(\alph*)]
  \item \emph{Creer que el problema es convexo sin comprobarlo}: un coste cóncavo a trozos (descuentos por volumen) o una restricción $x_1x_2\ge k$ rompen la convexidad; DCP lo detecta, un resolvedor local no.
  \item \emph{Signos de los duales}: cada herramienta tiene su convenio (\texttt{dual\_value} en CVXPY, \texttt{marginals} en SciPy); hay que comprobar el signo con una perturbación pequeña antes de usarlo como precio.
  \item \emph{Escalado}: mezclar MW con €/MWh$^2$ produce coeficientes de órdenes muy distintos y tolerancias que no significan lo mismo en todas las restricciones.
  \item \emph{Diferencias finitas en simuladores con ruido}: el gradiente numérico de una función ruidosa es basura; en ese caso conviene un método sin derivadas (\cref{tema:5}).
  \item \emph{Ignorar los límites de validez del modelo}: una curva de coste ajustada entre 100 y 400~MW no dice nada a 450~MW; los límites técnicos deben ser restricciones duras.
\end{enumerate}
\end{observacion}

\begin{lecturas}
\item Boyd y Vandenberghe~\cite{bv}, cap.~4 (§4.3 «Linear optimization problems», §4.4 «Quadratic optimization problems») y cap.~5 (§5.5 «Optimality conditions», §5.6 «Perturbation and sensitivity analysis»).
\item Nocedal y Wright~\cite{nw}, cap.~12 (teoría de la optimización con restricciones).
\item Wood, Wollenberg y Sheblé~\cite{wood}, cap.~3, despacho económico de grupos térmicos.
\item Documentación de CVXPY~\cite{cvxpydoc}: «Disciplined Convex Programming» y variables duales.
\end{lecturas}

\section*{Ejercicios del tema 4}
\addcontentsline{toc}{section}{Ejercicios del tema 4}

\begin{ejercicio}{1}{eje:4-dcp}
Indique si cada problema es convexo y por qué: (a) $\min x_1^{2}+x_2^{2}$ s.a. $x_1+x_2\ge 1$; (b) $\min -\log x_1-\log x_2$ s.a. $x_1+x_2=1$; (c) $\min x_1^{2}$ s.a. $x_1^{2}+x_2^{2}=1$; (d) $\min\ \max(x_1,x_2)$ s.a. $x\ge 0$, $x_1x_2\ge 1$; (e) $\min\ \norm{Ax-b}_2$ s.a. $\norm{x}_\infty\le 1$.
\end{ejercicio}

\begin{ejercicio}{2}{eje:4-kkt}
Resuelva con las condiciones KKT el problema $\min x_1^{2}+x_2^{2}$ s.a. $x_1+2x_2\ge 5$. Calcule el multiplicador e interprételo: ¿cuánto cambia aproximadamente el valor óptimo si el lado derecho pasa de 5 a $5{,}1$? Compruebe con el valor exacto.
\fuente{adaptado de Boyd y Vandenberghe~\cite{bv}, cap.~5, §5.5 (condiciones KKT) y §5.6 (sensibilidad).}
\end{ejercicio}

\begin{ejercicio}{2}{eje:4-despacho2}
Dos grupos tienen costes $C_1=20P_1+0{,}01P_1^{2}$ y $C_2=18P_2+0{,}02P_2^{2}$ (€/h) y deben cubrir 500~MW, sin límites activos. Obtenga el despacho por la regla del coste incremental igual y el precio marginal. Compruebe el resultado con el coste total.
\fuente{adaptado de Wood, Wollenberg y Sheblé~\cite{wood}, cap.~3, despacho económico por igualdad de costes incrementales.}
\end{ejercicio}

\begin{ejercicio}{2}{eje:4-agua}
Se reparte una potencia total de 1~W entre tres canales con niveles de ruido $\alpha=(0{,}5;\ 1;\ 1{,}5)$ para maximizar $\sum_i\log(\alpha_i+x_i)$ con $x\ge 0$ y $\sum_ix_i=1$. Escriba las KKT, deduzca que $x_i=\max(0,\ 1/\nu-\alpha_i)$ y calcule $\nu$ y la asignación.
\fuente{adaptado de Boyd y Vandenberghe~\cite{bv}, cap.~5, §5.5, ejemplo del llenado de agua (\emph{water-filling}).}
\end{ejercicio}

\begin{ejercicio}{2}{eje:4-sensib}
En el \cref{ej:despacho}: (a) estime con $\lambda$ el ahorro de reducir la demanda en 20~MW y compárelo con el valor exacto; (b) ¿hasta qué demanda sigue activo solo el límite superior del grupo~2, si los demás límites no cambian?
\end{ejercicio}

\begin{ejercicio}{2}{eje:4-dual}
Para el despacho sin límites técnicos del \cref{ej:despacho} (paso~1), calcule la función dual $q(\lambda)=\min_P\bigl[\sum_iC_i(P_i)-\lambda(\sum_iP_i-D)\bigr]$ de forma explícita y compruebe que su máximo se alcanza en $\lambda=28$ y vale lo mismo que el coste primal.
\end{ejercicio}

\begin{ejercicio}{3}{eje:4-primer-orden}
Sea $f$ convexa y diferenciable y $\mathcal{X}$ convexo. Demuestre que $x^{*}\in\mathcal{X}$ minimiza $f$ en $\mathcal{X}$ si y solo si $\nabla f(x^{*})\T(x-x^{*})\ge 0$ para todo $x\in\mathcal{X}$. Interprete la condición geométricamente y deduzca la condición $\nabla f(x^{*})=0$ cuando $x^{*}$ es interior.
\fuente{adaptado de Boyd y Vandenberghe~\cite{bv}, cap.~4, criterio de optimalidad para funciones diferenciables.}
\end{ejercicio}
